\section{Radar Systems and Passive Reception}\label{sec:bg:radar}

\subsection*{Pulsed radar}

The term \gls{radar} stands for \textbf{Ra}dio \textbf{D}etection \textbf{a}nd \textbf{R}anging.
A radar transmits a short electromagnetic pulse, waits for the echo that scatters from objects in the environment, and infers range and motion from that return \parencite{skolnik2008radar}.
The same principle supports air traffic control, weather observation, navigation, altitude measurement, and target tracking \parencite{skolnik2008radar}.

Distance follows from the round-trip delay $\tau$ between transmission and reception.
With wave speed $c$, the range of a monostatic radar---transmitter and receiver on the same platform---is
\begin{equation}
\label{eq:bg:range}
R = \frac{c\tau}{2},
\end{equation}
where the factor of 2 accounts for the outbound and return paths \parencite{skolnik2008radar}.
Animal echolocation uses the same timing argument with sound in place of radio waves \parencite{skolnik2008radar}.
Separated receive antennas, in bistatic or multistatic geometries, add angle of arrival and other spatial measurements to this range estimate.

Practical radars emit pulses at a pulse repetition interval (\gls{pri}), or equivalently a pulse repetition frequency (\gls{prf}).
If an echo from an earlier pulse arrives after the next pulse has been sent, the measured delay no longer identifies a unique range: this is \emph{range ambiguity}.
The largest range that remains unique is the \emph{unambiguous range}
\begin{equation}
\label{eq:bg:runamb}
R_{\mathrm{u}} = \frac{c}{2\,\mathrm{PRF}} = \frac{c\,\mathrm{PRI}}{2}.
\end{equation}
Matched filtering and coherent integration improve detection inside this window, but they cannot remove the ambiguity that pulse timing itself imposes \parencite{skolnik2008radar}.

\begin{figure}[tbp]
\centering
% =====================================================================
% Monostatic radar geometry and round-trip timing.
% A co-located TX/RX dish illuminates a target at range R; the transmitted
% pulse and its echo are paired with the corresponding events on a time
% axis to make the round-trip delay tau = 2R/c visible at a glance.
% Loaded from sections/02_background/01_radar_systems_passive_reception.tex
% via \input{...}.
% =====================================================================
\begin{tikzpicture}[
    >=Latex,
    font=\small,
    txstyle/.style={blue!60!black, thick},
    rxstyle/.style={orange!70!black, thick, dashed},
    annot/.style={font=\small\itshape},
    measure/.style={<->, thick, shorten >=1pt, shorten <=1pt},
  ]

  % --- geometry constants --------------------------------------------
  \def\Rrange{7.0}      % horizontal separation between radar and target
  \def\dishX{0.0}
  \def\targetX{\Rrange}

  % --- radar dish on the left ----------------------------------------
  \begin{scope}[shift={(\dishX,0)}]
    % parabolic reflector (filled crescent)
    \fill[blue!12] (-0.05,0.55) .. controls (-0.45,0) .. (-0.05,-0.55)
      -- (0.18,-0.55) .. controls (-0.22,0) .. (0.18,0.55) -- cycle;
    \draw[blue!60!black, thick]
      (-0.05,0.55) .. controls (-0.45,0) .. (-0.05,-0.55);
    \draw[blue!60!black, thick]
      (0.18,0.55) .. controls (-0.22,0) .. (0.18,-0.55);
    \draw[blue!60!black, thick] (-0.05,0.55) -- (0.18,0.55);
    \draw[blue!60!black, thick] (-0.05,-0.55) -- (0.18,-0.55);
    % feed and stand
    \draw[thick] (0.07,-0.55) -- (0.07,-1.0);
    \draw[thick] (-0.25,-1.0) -- (0.40,-1.0);
    \node[align=center, below, font=\footnotesize] at (0.07,-1.05)
      {Radar\\co-located TX/RX};
  \end{scope}

  % --- aircraft silhouette on the right ------------------------------
  \begin{scope}[shift={(\targetX,0)}]
    % fuselage
    \draw[thick, fill=gray!25] (-0.55,-0.10) -- (0.55,-0.10)
      -- (0.65,0.02) -- (0.55,0.15) -- (-0.50,0.15) -- (-0.60,0.02) -- cycle;
    % main wings
    \draw[thick, fill=gray!35] (-0.05,0.15) -- (-0.45,0.55)
      -- (0.20,0.55) -- (0.30,0.15) -- cycle;
    \draw[thick, fill=gray!35] (-0.05,-0.10) -- (-0.45,-0.50)
      -- (0.20,-0.50) -- (0.30,-0.10) -- cycle;
    % vertical stabilizer
    \draw[thick, fill=gray!35] (0.40,0.15) -- (0.35,0.45)
      -- (0.60,0.45) -- (0.60,0.15) -- cycle;
    \node[above, font=\footnotesize] at (0,0.65) {Target};
  \end{scope}

  % --- transmitted pulse (radar -> target) ---------------------------
  \draw[txstyle, ->] (0.4,0.10) -- (\targetX - 0.65,0.10);
  \node[annot, blue!60!black, above]
    at ({(\dishX + \targetX)/2}, 0.10) {transmitted pulse};

  % --- echo (target -> radar) ----------------------------------------
  \draw[rxstyle, ->] (\targetX - 0.65,-0.10) -- (0.4,-0.10);
  \node[annot, orange!70!black, below]
    at ({(\dishX + \targetX)/2}, -0.10) {echo};

  % --- range arrow ---------------------------------------------------
  \draw[measure] (\dishX + 0.07,-1.45) -- (\targetX,-1.45);
  \node[annot, fill=white, inner sep=2pt]
    at ({(\dishX + \targetX)/2}, -1.45) {range $R$};

  % --- timeline below the scene --------------------------------------
  \def\tlY{-2.6}
  \def\txT{0.5}
  \def\rxT{5.5}

  \draw[->] (-0.2,\tlY) -- (\Rrange + 0.5,\tlY) node[right] {time $t$};
  \draw[txstyle, line width=1pt] (\txT,\tlY) -- (\txT,\tlY + 0.7);
  \node[annot, blue!60!black, above] at (\txT,\tlY + 0.7) {TX};
  \draw[rxstyle, line width=1pt] (\rxT,\tlY) -- (\rxT,\tlY + 0.55);
  \node[annot, orange!70!black, above] at (\rxT,\tlY + 0.55) {RX (echo)};
  \draw[measure] (\txT,\tlY - 0.45) -- (\rxT,\tlY - 0.45);
  \node[annot, fill=white, inner sep=2pt]
    at ({(\txT + \rxT)/2}, \tlY - 0.45) {$\tau = 2R/c$};

\end{tikzpicture}

\caption{Monostatic radar geometry and round-trip timing. A co-located transmit/receive (TX/RX) aperture on the left emits a pulse that travels to a target at range $R$, scatters, and returns to the same aperture as an echo. The lower timeline shows the corresponding events: the outgoing TX pulse at $t=0$ and the received echo at $t=\tau$, separated by the round-trip delay $\tau = 2R/c$ from \cref{eq:bg:range}.}
\label{fig:bg:pulse_echo_timeline}
\end{figure}

\subsection*{Active radar versus passive reception}

A transmitting radar measures delay against its own outgoing pulses and thereby obtains range \parencite{skolnik2008radar}.
Those same pulses are visible to any receiver in range, so transmission reveals presence, location, and activity as well as it measures them \parencite{wiley2006elint}.

Passive systems such as \gls{esm} and \gls{elint} do not transmit.
They listen to other radars to determine which emitters are present and how those emitters are behaving \parencite{wiley2006elint}.
Without control of the waveform, the receiver is limited to what it can measure after detection: carrier frequency, time of arrival, pulse width, amplitude, and, if the antenna supports it, angle of arrival.
Those measurements are stored as the \gls{pdw} defined next.
How a classical processing chain uses them is the subject of \cref{sec:bg:elint_pipeline}.

\subsection{Pulse descriptor words}\label{sec:bg:pdw}

A receiver records a voltage time series at the antenna.
Pulses appear as brief energy rises above the noise floor.
For each detection the receiver estimates carrier frequency, pulse width, time of arrival, amplitude, and, when directional antennas are available, angle of arrival.

Those estimates are stored as a \emph{\gls{pdw}}:
\begin{equation}\label{eq:bg:pdw_vector}
  \vect{x}_n
  =
  \bigl(
    f_n,\,
    \mathrm{PW}_n,\,
    t_n,\,
    A_n,\,
    \theta_n,\,
    \ldots
  \bigr)^{\top}
  \in \R^{d},
\end{equation}
where $f_n$ is the measured carrier frequency, $\mathrm{PW}_n$ is pulse width, $t_n$ is time of arrival (often the leading edge), $A_n$ is amplitude, and $\theta_n$ is angle of arrival if available.
The ellipsis allows further recorded properties, such as elevation; once the template is chosen, every pulse is stored in the same format.
\cref{fig:bg:pdw} shows this five-parameter template.

\begin{figure}[tbp]
  \centering
  % =====================================================================
% Schematic of the parameters carried by a pulse descriptor word (PDW)
% for use in \cref{sec:bg:pdw}.  A short pulse train illustrates the
% time-domain PDW fields (TOA, PW, PRI, amplitude, RF carrier); a small
% inset shows AOA as a spatial bearing on the receive aperture.
% Loaded from sections/02_background/01_radar_systems_passive_reception.tex via \input{...}.
% =====================================================================
\begin{tikzpicture}[
    >=Latex,
    font=\small,
    pulse/.style={draw=blue!60!black, thick, fill=blue!10},
    rfwave/.style={very thin, red!60!black},
    annot/.style={font=\small\itshape},
    measure/.style={<->, thick, shorten >=1pt, shorten <=1pt},
    drop/.style={dashed, thin, gray!70!black},
  ]

  % --- pulse train geometry (TOA = leading edge, all PW equal here) ---
  \def\toaA{1.4}
  \def\toaB{4.4}
  \def\toaC{7.4}
  \def\pw{0.9}
  \def\ampA{2.4}
  \def\ampB{2.2}
  \def\ampC{2.5}

  % --- axes -----------------------------------------------------------
  \draw[->] (-0.2, 0) -- (10.4, 0) node[below] {time $t$};
  \draw[->] (0, -0.2) -- (0, 3.6) node[left]  {amplitude};

  % --- pulse envelopes ------------------------------------------------
  \draw[pulse] (\toaA, 0) rectangle (\toaA + \pw, \ampA);
  \draw[pulse] (\toaB, 0) rectangle (\toaB + \pw, \ampB);
  \draw[pulse] (\toaC, 0) rectangle (\toaC + \pw, \ampC);

  % --- illustrative RF carrier inside the middle pulse ---------------
  \draw[rfwave, samples=120, smooth, domain=\toaB:\toaB+\pw]
    plot (\x, {0.5*\ampB + 0.42*\ampB*sin(2880*(\x - \toaB))});

  % --- TOA drop lines and labels --------------------------------------
  \foreach \t/\lbl in {%
      \toaA/$t_n$,%
      \toaB/$t_{n+1}$,%
      \toaC/$t_{n+2}$%
    } {%
    \draw[drop] (\t, 0) -- (\t, -0.35);
    \node[annot, below] at (\t, -0.32) {\lbl};
  }
  \node[annot, below right=1pt and -2pt] at (\toaA, -0.32)
    {\scriptsize\hspace{0.15em}(TOA)};

  % --- amplitude arrow on the first pulse -----------------------------
  \draw[measure] (\toaA - 0.5, 0) -- (\toaA - 0.5, \ampA);
  \node[annot, left] at (\toaA - 0.5, \ampA/2) {$A_n$};

  % --- PW arrow above the third pulse ---------------------------------
  \draw[measure] (\toaC, \ampC + 0.3) -- (\toaC + \pw, \ampC + 0.3);
  \node[annot, above] at (\toaC + \pw/2, \ampC + 0.3) {PW};

  % --- PRI arrow between the first two leading edges ------------------
  \pgfmathsetmacro{\priY}{\ampA + 0.95}
  \draw[measure] (\toaA, \priY) -- (\toaB, \priY);
  \node[annot, above] at ({(\toaA + \toaB)/2}, \priY) {PRI $= 1/$PRF};

  % --- RF label with leader line into the middle-pulse oscillation ----
  \node[annot, text=red!60!black] (rflabel) at (6.4, 0.55*\ampB)
    {RF $f_n$};
  \draw[->, very thin, red!60!black]
    (rflabel.west) -- (\toaB + \pw - 0.05, 0.55*\ampB);

  % --- AOA inset (upper right): bearing on the receive aperture -------
  \begin{scope}[shift={(9.0, 2.55)}]
    % receive aperture
    \draw[thick] (-0.30, 0) -- (0.30, 0) -- (0, 0.45) -- cycle;
    \node[font=\footnotesize, below=1pt] at (0, 0) {Rx};
    % boresight reference
    \draw[drop] (0, 0.45) -- (0, 1.55);
    % incoming wavefront direction
    \draw[->, thick] (1.0, 1.40) -- (0.04, 0.50);
    % angle arc
    \draw[->, thin] (0, 1.15) arc (90:50:0.7);
    \node[annot] at (0.55, 1.00) {AOA $\theta_n$};
  \end{scope}

\end{tikzpicture}

  \caption{Schematic of the parameters captured by a \gls{pdw} for a short pulse train. Each pulse contributes a \gls{toa} (the leading edge $t_n$ on the time axis), a pulse width (PW), a peak amplitude $A_n$, and a carrier \gls{rf} $f_n$, sketched as the fast oscillation inside the middle pulse; successive pulses also define a \gls{pri} (equivalently a \gls{prf} $=1/\mathrm{PRI}$). The inset shows the \gls{aoa} $\theta_n$ measured at the receive aperture relative to a boresight reference. A typical \gls{pdw} stacks these per-pulse measurements into a vector $\vect{x}_n = (f_n, \mathrm{PW}_n, t_n, A_n, \theta_n)$ that is the usual input to downstream emitter and mode reasoning \parencite{wiley2006elint}.}
  \label{fig:bg:pdw}
\end{figure}

A sequence of pulses from one transmitter is a \emph{pulse train}.
When that transmitter holds a stable mode, the inter-arrival times
\begin{equation}\label{eq:bg:pdw_pri}
  \tau_n = t_n - t_{n-1}, \qquad n \ge 2
\end{equation}
are regular or follow a programmed schedule; that interval is the \gls{pri} of the train, and its inverse is the \gls{prf}.
Simple emitters keep a fixed interval; agile emitters vary it in a planned way.

Each measured descriptor is a noisy, possibly incomplete observation of the true pulse.
Weak pulses may be missed, and noise events may be reported as pulses.

After this step the sampled voltage is discarded.
The working input is a time-ordered sequence of \glspl{pdw} $(\vect{x}_1,\ldots,\vect{x}_T)$, called a \emph{recording} because it generally interleaves the pulse trains of several transmitters.

\subsection{Emitters and operating modes}\label{sec:bg:emitters}

In passive \gls{esm} and \gls{elint}, an \emph{emitter} is a physical transmitter: pulses that line up in time, frequency, and direction strongly enough to be treated as one source \parencite{wiley2006elint}.
An \emph{operating mode} is the transmit schedule in force at a given time---carrier plan, pulse width, repetition pattern, and scan or switching rules---chosen for a task such as search, acquisition, track, or illumination \parencite{skolnik2008radar,wiley2006elint}.

These are two different groupings of the same pulses.
One emitter switches among modes as its task changes, so a single physical source produces several mode segments over a mission.
Different emitters can also run the same functional mode with different numerical parameters, for example two trackers with distinct nominal \glspl{prf} \parencite{wiley2006elint}.

Operators therefore care about both questions: which transmitter is present, and what that transmitter is doing now.
\Cref{sec:bg:elint_pipeline} describes how classical processing answers them in sequence.
